\documentclass[11pt,a4paper]{article}

% --- Codificación e idioma ---
\usepackage[utf8]{inputenc}
\usepackage[T1]{fontenc}
\usepackage[spanish,es-tabla,es-noshorthands]{babel}
\usepackage{lmodern}

% --- Geometría y formato general ---
\usepackage[margin=2.5cm]{geometry}
\usepackage{microtype}
\usepackage{parskip}
\usepackage{enumitem}
\setlist{topsep=2pt,itemsep=2pt,parsep=0pt}

% --- Tablas y diagramas ---
\usepackage{array}
\usepackage{booktabs}
\usepackage{amsmath,amssymb}
\usepackage{tikz}
\usetikzlibrary{positioning,arrows.meta,calc,shapes.geometric}

% --- Cabeceras y numeración ---
\usepackage{fancyhdr}
\pagestyle{fancy}
\fancyhf{}
\fancyhead[L]{\small Bases de Datos II}
\fancyhead[R]{\small Resumen --- Transacciones y Concurrencia}
\fancyfoot[C]{\thepage}
\renewcommand{\headrulewidth}{0.4pt}

% --- Color y código ---
\usepackage{xcolor}
\definecolor{codebg}{RGB}{248,248,248}
\definecolor{codeframe}{RGB}{200,200,200}
\definecolor{codekw}{RGB}{0,0,180}
\definecolor{codestr}{RGB}{163,21,21}
\definecolor{codecmt}{RGB}{0,128,0}
\definecolor{notebg}{RGB}{255,250,230}
\definecolor{noteframe}{RGB}{220,180,80}

\usepackage{listings}
\lstdefinelanguage{SQLext}[]{SQL}{
    morekeywords={BEGIN,COMMIT,ROLLBACK,START,TRANSACTION,
        ISOLATION,LEVEL,SESSION,
        SERIALIZABLE,REPEATABLE,READ,COMMITTED,UNCOMMITTED,
        NOWAIT,REFERENCES,
        DEFERRABLE,INITIALLY,DEFERRED,IMMEDIATE,
        ALTER,CONSTRAINT,CONSTRAINTS},
    sensitive=false,
    morecomment=[l]{--},
    morestring=[b]'
}
\lstdefinestyle{sqlstyle}{
    language=SQLext,
    basicstyle=\ttfamily\footnotesize,
    keywordstyle=\color{codekw}\bfseries,
    stringstyle=\color{codestr},
    commentstyle=\color{codecmt}\itshape,
    backgroundcolor=\color{codebg},
    frame=single,
    rulecolor=\color{codeframe},
    breaklines=true,
    columns=fullflexible,
    keepspaces=true,
    showstringspaces=false,
    captionpos=b,
    xleftmargin=0.4em,
    xrightmargin=0.4em,
    aboveskip=8pt,
    belowskip=8pt,
    tabsize=2
}
\lstset{style=sqlstyle}

% --- Cajas de nota ---
\usepackage[most]{tcolorbox}
\newtcolorbox{nota}{colback=notebg,colframe=noteframe,arc=2pt,
    left=6pt,right=6pt,top=4pt,bottom=4pt,boxrule=0.6pt}

% --- Hipervínculos ---
\usepackage[hidelinks,bookmarks=true]{hyperref}

\title{Resumen Teórico --- Transacciones y Control de Concurrencia\\[4pt]
       \large Bases de Datos II --- Práctico N.\textsuperscript{o} 6}
\author{Tomás Rodeghiero}
\date{2026}

\begin{document}

\begin{titlepage}
    \centering
    \vspace*{1.5cm}

    {\large \textbf{Universidad Nacional de Río Cuarto}}\\[6pt]
    {\normalsize Facultad de Ciencias Exactas, Físico-Químicas y Naturales}\\[4pt]
    {\normalsize Departamento de Computación}\\[4pt]
    {\normalsize Licenciatura en Ciencias de la Computación}\\[3cm]

    {\Large \textbf{Bases de Datos II}}\\[1.2cm]

    {\Large \textbf{Resumen Teórico}}\\[6pt]
    {\large Transacciones y Control de Concurrencia\\[2pt]
    Práctico N.\textsuperscript{o} 6}\\[3cm]

    \begin{tabular}{rl}
        \textbf{Alumno:} & Tomás Rodeghiero \\[4pt]
        \textbf{Año:}    & 2026 \\
    \end{tabular}

    \vfill
    {\small Síntesis del Teórico 8: Transacciones y Control de Concurrencia.}
\end{titlepage}

\tableofcontents
\newpage

% =====================================================
\section{Qué es una transacción}
% =====================================================

Una \textbf{transacción} es una unidad lógica de trabajo de una aplicación
que accede ---y posiblemente modifica--- elementos de la base de datos.
Para el motor, una transacción es una secuencia de operaciones
$r(Q)$ (lectura de un ítem $Q$) y $w(Q)$ (escritura) terminada por una
operación de finalización: \texttt{COMMIT} si se completó con éxito o
\texttt{ROLLBACK} si hubo que abortarla.

\paragraph{La regla básica.}
\emph{La base de datos parte de un estado consistente y, después de
ejecutar una transacción, debe quedar consistente otra vez.} Durante la
ejecución la base puede pasar temporalmente por estados inconsistentes
---por ejemplo, A ya descontó $\$500$ pero B todavía no los recibió---,
pero esos estados son privados de la transacción: nadie de afuera debería
verlos.

\paragraph{Ejemplo canónico.} Transferir $\$500$ de la cuenta $A$ a la $B$:
\[
\texttt{read}(A);\; A := A - 500;\; \texttt{write}(A);\;
\texttt{read}(B);\; B := B + 500;\; \texttt{write}(B).
\]
Las garantías que se le piden a la transacción son tres:

\begin{itemize}
    \item Si el sistema falla entre el \texttt{write}(A) y el
    \texttt{write}(B), el cambio sobre $A$ \emph{no} debe quedar
    reflejado: la base quedaría inconsistente.
    \item Ninguna otra transacción debe ver el estado intermedio (con $A$
    ya descontado y $B$ aún sin acreditar).
    \item Una vez que la transacción confirma, los $\$500$ están
    transferidos y eso se mantiene aunque después falle hardware o
    software.
\end{itemize}

% =====================================================
\section{Propiedades ACID}
% =====================================================

Las propiedades que un motor debe asegurar a cada transacción se
resumen en el acrónimo \textbf{ACID}:

\begin{itemize}
    \item \textbf{Atomicidad (A).} O todas las operaciones de la transacción
    se realizan, o ninguna. No hay ``a medias''. La implementa el manejador
    de recuperación (logs + rollback).
    \item \textbf{Consistencia (C).} La ejecución aislada de una transacción
    \emph{correctamente programada} preserva la consistencia de la base.
    Esta propiedad cae en la responsabilidad del programador: el motor
    no sabe qué reglas de negocio son ``consistentes''.
    \item \textbf{Aislamiento (I).} Aunque varias transacciones se ejecuten
    concurrentemente, cada una percibe la base como si fuera la única que
    está actuando sobre ella. Es lo que implementa el control de
    concurrencia (bloqueos, marcas temporales, MVCC).
    \item \textbf{Durabilidad (D).} Una vez que una transacción confirmó
    (\texttt{COMMIT}), sus cambios persisten en disco aun ante fallos
    posteriores. Lo asegura el sistema de logs + flushes.
\end{itemize}

% =====================================================
\section{Estados de una transacción}
% =====================================================

Durante su vida, una transacción pasa por una secuencia de estados:

\begin{itemize}
    \item \textbf{Activa}: está ejecutándose normalmente.
    \item \textbf{Parcialmente comprometida}: ejecutó su última instrucción
    pero todavía hay trabajo del motor pendiente (escribir log, propagar
    cambios).
    \item \textbf{Fallida}: algo impide continuar la ejecución (error en
    la aplicación, violación de constraint, deadlock, fallo del sistema).
    \item \textbf{Abortada}: se hizo rollback y la base se restauró al
    estado previo. Una transacción abortada se puede \emph{reintentar}
    (si el error fue transitorio, como un deadlock) o \emph{cancelar
    definitivamente}.
    \item \textbf{Comprometida}: se completó con éxito; sus efectos son
    permanentes.
\end{itemize}

\begin{center}
\begin{tikzpicture}[
    node distance=14mm and 22mm,
    estado/.style={draw,ellipse,font=\footnotesize,inner sep=3pt}
]
    \node[estado] (a)  {Activa};
    \node[estado] (pc) [above right=of a]   {Parcialm. comprom.};
    \node[estado] (c)  [right=of pc]        {Comprometida};
    \node[estado] (f)  [below right=of a]   {Fallida};
    \node[estado] (ab) [right=of f]         {Abortada};

    \draw[->] (a)  -- (pc);
    \draw[->] (pc) -- (c);
    \draw[->] (a)  -- (f);
    \draw[->] (pc) -- (f);
    \draw[->] (f)  -- (ab);
\end{tikzpicture}
\end{center}

% =====================================================
\section{Ejecución concurrente y planificaciones}
% =====================================================

Permitir que muchas transacciones corran al mismo tiempo trae beneficios
concretos:

\begin{itemize}
    \item \textbf{Mejor uso del hardware}: mientras una transacción usa
    CPU, otra puede estar leyendo disco.
    \item \textbf{Menor tiempo de respuesta promedio}: transacciones
    cortas no tienen que esperar detrás de transacciones largas.
\end{itemize}

Pero también introduce problemas: si las transacciones se entrelazan
sin control, pueden violar la consistencia. Por eso hace falta un
mecanismo de \textbf{control de concurrencia} que regule cómo se
intercalan.

\subsection{Planificación (\emph{schedule})}

Una \textbf{planificación} es una secuencia que muestra en qué orden se
ejecutan las operaciones de un conjunto de transacciones concurrentes.
Tiene dos requisitos:

\begin{itemize}
    \item Contiene todas las operaciones de todas las transacciones.
    \item Preserva el orden interno de cada transacción.
\end{itemize}

\paragraph{Planificación serial.} Ejecuta cada transacción completa antes
de empezar la siguiente; no hay intercalado. Son obviamente correctas
porque cada transacción ve la base entera entre las dos consistentes.
La idea de la concurrencia es \emph{aproximarse} a una ejecución serial
en términos de resultado, pero permitiendo intercalado para ganar
rendimiento.

% =====================================================
\section{Serializabilidad}
% =====================================================

Asumimos dos cosas:
\begin{itemize}
    \item Cada transacción tomada aisladamente preserva la consistencia
    de la base.
    \item Cualquier planificación serial preserva la consistencia.
\end{itemize}

Una planificación (posiblemente concurrente) es \textbf{secuenciable}
(o \emph{serializable}) si es \emph{equivalente} a alguna planificación
serial. Hay dos definiciones de equivalencia:

\begin{itemize}
    \item Equivalencia \textbf{por conflicto} (la más usada y la del
    Práctico~6).
    \item Equivalencia \textbf{por vistas} (más permisiva pero NP-hard).
\end{itemize}

\subsection{Conflicto entre operaciones}

Dos instrucciones $I_i$ de $T_i$ y $I_j$ de $T_j$ están en conflicto si
acceden al mismo ítem $Q$ y al menos una de las dos es una escritura:

\begin{center}
\begin{tabular}{lll}
\toprule
$I_i$ & $I_j$ & ¿Conflicto? \\
\midrule
$r(Q)$ & $r(Q)$ & No \\
$r(Q)$ & $w(Q)$ & Sí \\
$w(Q)$ & $r(Q)$ & Sí \\
$w(Q)$ & $w(Q)$ & Sí \\
\bottomrule
\end{tabular}
\end{center}

\subsection{Equivalencia por conflicto}

Dos planificaciones son \textbf{equivalentes por conflicto} si una se
puede transformar en la otra intercambiando \emph{únicamente}
instrucciones consecutivas que no estén en conflicto. Una planificación
es \textbf{serializable por conflicto} si es equivalente por conflicto
a alguna planificación serial.

\subsection{Test del grafo de precedencia}

El test algorítmico de serializabilidad por conflicto construye un grafo
dirigido:

\begin{itemize}
    \item Un nodo por cada transacción de la planificación.
    \item Un arco $T_i \rightarrow T_j$ por cada par de operaciones en
    conflicto donde la de $T_i$ aparece \emph{antes} que la de $T_j$.
\end{itemize}

\begin{nota}
La planificación es serializable por conflicto si y sólo si el grafo
de precedencia es \textbf{acíclico}. Si hay un ciclo, ningún orden serial
es compatible con todos los arcos a la vez, y por lo tanto no hay una
serie equivalente.
\end{nota}

\paragraph{Ejemplo (Práctico 6, Ejercicio 1).}
La planificación
$r_1(X),\, r_2(X),\, w_1(X),\, w_2(X),\ldots$
genera $T_1 \rightarrow T_2$ (por $r_1(X)$ antes de $w_2(X)$) y
$T_2 \rightarrow T_1$ (por $r_2(X)$ antes de $w_1(X)$): hay ciclo,
y la planificación no es serializable.

% =====================================================
\section{Recuperabilidad}
% =====================================================

La serializabilidad asegura que el resultado de la concurrencia
``equivale'' a una ejecución serial. Pero falta otra dimensión: ¿qué pasa
si una transacción aborta? Para que se pueda hacer \texttt{ROLLBACK}
correctamente, la planificación tiene que cumplir condiciones extra.

\subsection{Planificación recuperable}

Una planificación es \textbf{recuperable} si, para todo par $T_i, T_j$
tal que $T_j$ lee un dato que $T_i$ escribió antes, el \texttt{commit}
de $T_i$ aparece antes que el de $T_j$. Si no se cumple, $T_j$ podría
haber confirmado basándose en una escritura que después se aborta:
\emph{no es posible deshacer eso}.

\subsection{Planificación sin cascada}

Una planificación es \textbf{sin cascada} (\emph{cascadeless}) si para
todo $T_i, T_j$ tal que $T_j$ lee un dato escrito por $T_i$, el
\texttt{commit} de $T_i$ ocurre antes que la \emph{lectura} de $T_j$.

\begin{itemize}
    \item Toda planificación sin cascada es recuperable.
    \item Lo recíproco no se cumple: hay recuperables que sí tienen
    cascada.
\end{itemize}

El \textbf{rollback en cascada} (que una falla obligue a abortar muchas
transacciones encadenadas) es indeseable porque desperdicia trabajo y
puede crear avalanchas. SQL exige que el sistema garantice tanto
secuencialidad como ausencia de cascada.

% =====================================================
\section{Niveles de aislamiento (SQL\,92)}
% =====================================================

Garantizar \emph{serializabilidad estricta} para todas las transacciones
es caro (más bloqueos, menos concurrencia). Para algunas aplicaciones
---típicamente reportes y estadísticas--- se acepta perder precisión a
cambio de rendimiento. SQL\,92 define cuatro niveles de aislamiento
ordenados de más débil a más fuerte:

\begin{center}
\begin{tabular}{lccc}
\toprule
Nivel & Lectura sucia & No repetible & Phantom \\
\midrule
\texttt{READ UNCOMMITTED} & Posible    & Posible    & Posible \\
\texttt{READ COMMITTED}   & No posible & Posible    & Posible \\
\texttt{REPEATABLE READ}  & No posible & No posible & Posible (raro) \\
\texttt{SERIALIZABLE}     & No posible & No posible & No posible \\
\bottomrule
\end{tabular}
\end{center}

\paragraph{Los tres fenómenos.}
\begin{itemize}
    \item \textbf{Lectura sucia (\emph{dirty read})}: $T_1$ lee datos que
    $T_2$ todavía no confirmó. Si $T_2$ aborta, $T_1$ leyó un fantasma.
    \item \textbf{Lectura no repetible}: $T_1$ lee la misma fila dos
    veces y obtiene valores distintos porque entre lectura y lectura
    otra transacción la actualizó y confirmó.
    \item \textbf{Phantom read}: $T_1$ ejecuta dos veces la misma
    consulta de rango y aparecen filas nuevas porque otra transacción
    insertó registros que cumplen la condición.
\end{itemize}

\paragraph{Comandos SQL.}
\begin{lstlisting}
-- MySQL
SET SESSION TRANSACTION ISOLATION LEVEL SERIALIZABLE;

-- PostgreSQL
SET DEFAULT_TRANSACTION_ISOLATION = 'serializable';
\end{lstlisting}

% =====================================================
\section{Control de concurrencia: protocolos basados en bloqueos}
% =====================================================

Los protocolos de bloqueo garantizan aislamiento haciendo que cada
operación pase por un \emph{lock manager} antes de poder ejecutarse.

\subsection{Modos de bloqueo}

\begin{itemize}
    \item \textbf{Compartido (S/C)}: \texttt{lock-S(Q)}. El dato $Q$ puede
    ser leído. Varias transacciones pueden tener S sobre $Q$ a la vez.
    \item \textbf{Exclusivo (X)}: \texttt{lock-X(Q)}. El dato $Q$ puede
    ser leído y escrito. Sólo una transacción puede tenerlo a la vez.
\end{itemize}

Matriz de compatibilidad:

\begin{center}
\begin{tabular}{c|cc}
\toprule
       & S         & X         \\
\midrule
S      & compatible & no compatible \\
X      & no compatible & no compatible \\
\bottomrule
\end{tabular}
\end{center}

Si una transacción $T_i$ pide un bloqueo que choca con uno ya tomado,
$T_i$ \textbf{espera} hasta que se libere.

\begin{nota}
\textbf{Bloquear no alcanza por sí solo.} Si una transacción pide y
suelta locks alegremente, distintos ítems pueden ser leídos de versiones
inconsistentes entre sí, y el resultado no equivale a ninguna ejecución
serial. Se necesita un \emph{protocolo} que diga cuándo se pueden tomar
y liberar.
\end{nota}

\subsection{Posibilidad de Deadlock}

Si $T_1$ tiene $X(B)$ y pide $X(A)$, y $T_2$ tiene $S(A)$ y pide $S(B)$
\emph{(o $X(B)$)}, ninguna avanza. Esto es un \textbf{deadlock}. Para
desbloquear, el motor debe hacer rollback de al menos una de las dos
(la \emph{víctima}); sus locks se liberan y la otra continúa.

% =====================================================
\section{Protocolo de bloqueo de dos fases (2PL)}
% =====================================================

El protocolo más usado para garantizar serializabilidad por conflicto
mediante bloqueos es el \textbf{Bloqueo de Dos Fases}:

\begin{itemize}
    \item \textbf{Fase de crecimiento}: la transacción puede solicitar
    bloqueos. No puede liberar ninguno.
    \item \textbf{Fase de decrecimiento}: la transacción puede liberar
    bloqueos. No puede solicitar ninguno nuevo.
\end{itemize}

El punto donde una transacción libera su primer lock se llama
\emph{punto de bloqueo}. A partir de ahí, todo lo que puede hacer es
operar con lo que ya tiene y soltar locks.

\paragraph{Garantía.} Toda planificación legal bajo 2PL es serializable
por conflicto. La prueba se basa en mostrar que el orden de los puntos
de bloqueo de las transacciones induce un orden serial equivalente.

\subsection{Variantes}

\begin{itemize}
    \item \textbf{2PL puro o básico}: el descrito arriba. Permite
    rollback en cascada (porque una transacción puede haber leído de
    otra que después aborta).
    \item \textbf{2PL estricto}: todos los \emph{bloqueos exclusivos} se
    mantienen hasta el final (commit/rollback). Evita el rollback en
    cascada.
    \item \textbf{2PL riguroso}: \emph{todos} los bloqueos (S y X) se
    retienen hasta el final. Asegura sin cascada y, además, las
    transacciones se serializan en el orden de sus commits. Es la
    variante que implementan la mayoría de los motores comerciales.
\end{itemize}

\subsection{Conversión de bloqueos}

Para ganar concurrencia, 2PL puede admitir \emph{conversiones}:

\begin{itemize}
    \item Subir: de S a X (sólo durante la fase de crecimiento).
    \item Bajar: de X a S (sólo durante la fase de decrecimiento).
\end{itemize}

Si los locks se retienen hasta el final, las conversiones son seguras y
las planificaciones quedan sin cascada.

% =====================================================
\section{Protocolos basados en grafos}
% =====================================================

Otra alternativa al 2PL son los protocolos basados en grafos. Imponen
un orden parcial $\rightarrow$ sobre los ítems de datos: si
$d_i \rightarrow d_j$, toda transacción que use ambos debe bloquear
$d_i$ \emph{antes} de $d_j$.

\subsection{Protocolo de árbol}

El conjunto de ítems se organiza como un árbol enraizado. Sólo permite
bloqueos exclusivos. Reglas:

\begin{enumerate}
    \item El \emph{primer} bloqueo de $T_i$ puede ser sobre cualquier ítem.
    \item Los siguientes ítems sólo pueden bloquearse si $T_i$ tiene
    actualmente bloqueado el padre.
    \item Los ítems se pueden desbloquear en cualquier momento.
    \item $T_i$ no puede volver a bloquear un ítem que ya bloqueó y
    desbloqueó.
\end{enumerate}

\paragraph{Ventajas frente a 2PL.}
\begin{itemize}
    \item No hay deadlocks (porque el orden está fijo de antemano).
    \item Las transacciones pueden liberar locks tempranamente, ganando
    concurrencia.
\end{itemize}

\paragraph{Desventajas.}
\begin{itemize}
    \item Una transacción puede tener que bloquear ítems que no necesita,
    sólo para llegar a los que sí (\emph{lock pre-emption}).
    \item Las planificaciones pueden no ser recuperables ni sin cascada.
\end{itemize}

% =====================================================
\section{Protocolo de ordenación por marcas temporales}
% =====================================================

La alternativa a los bloqueos es decidir \emph{a priori} el orden serial
equivalente y forzar al schedule a respetarlo. Para eso cada transacción
$T_i$ recibe al iniciar una \textbf{marca temporal} $MT(T_i)$ única
(reloj del sistema o contador). Si $MT(T_i) < MT(T_j)$, el motor se
compromete a producir un schedule equivalente a la serie en la que $T_i$
va antes que $T_j$.

\paragraph{Información que mantiene cada ítem.}
\begin{itemize}
    \item $\mathrm{MT\text{-}E}(Q)$: mayor marca de las transacciones que
    escribieron $Q$ con éxito.
    \item $\mathrm{MT\text{-}L}(Q)$: mayor marca de las transacciones que
    leyeron $Q$ con éxito.
\end{itemize}

\subsection{Reglas del protocolo}

\paragraph{$T_i$ ejecuta $\mathrm{read}(Q)$.}
\begin{itemize}
    \item Si $MT(T_i) < \mathrm{MT\text{-}E}(Q)$: $T_i$ querría leer una
    versión que ya fue sobrescrita por una transacción más nueva. La
    lectura se rechaza y $T_i$ se retrocede.
    \item En otro caso: se ejecuta y
    $\mathrm{MT\text{-}L}(Q) \gets \max(\mathrm{MT\text{-}L}(Q), MT(T_i))$.
\end{itemize}

\paragraph{$T_i$ ejecuta $\mathrm{write}(Q)$.}
\begin{itemize}
    \item Si $MT(T_i) < \mathrm{MT\text{-}L}(Q)$: el valor que $T_i$
    produciría ya fue necesario por una lectura posterior; la escritura
    llega ``tarde'' y $T_i$ se retrocede.
    \item Si $MT(T_i) < \mathrm{MT\text{-}E}(Q)$: $T_i$ querría
    sobrescribir con un valor obsoleto; se retrocede.
    \item En otro caso: se ejecuta y
    $\mathrm{MT\text{-}E}(Q) \gets MT(T_i)$.
\end{itemize}

\subsection{Propiedades}

\begin{itemize}
    \item Garantiza planificaciones serializables.
    \item No tiene deadlocks (las transacciones nunca esperan; se retroceden).
    \item Puede generar planificaciones \emph{no recuperables}. Para
    asegurar recuperabilidad hay que extenderlo (por ejemplo, retrasar
    el \texttt{commit} de $T_i$ hasta que todas las que escribieron datos
    leídos por $T_i$ hayan commiteado).
    \item Costo: las transacciones largas suelen ser retrocedidas y
    reiniciadas (\emph{starvation}).
\end{itemize}

% =====================================================
\section{Esquemas multiversión y MVCC}
% =====================================================

Los esquemas \textbf{multiversión} mantienen varias versiones de cada
ítem de datos para incrementar la concurrencia. Cuando una transacción
lee, el control de concurrencia le entrega la versión adecuada según su
marca temporal o su snapshot, sin tener que bloquear escritores.

\subsection{Ordenación por marcas temporales multiversión}

Cada ítem $Q$ tiene una secuencia de versiones $Q_1, Q_2, \ldots, Q_m$.
Cada versión $Q_k$ guarda:

\begin{itemize}
    \item \texttt{contenido}: el valor de la versión.
    \item $\mathrm{MT\text{-}E}(Q_k)$: la marca temporal de la
    transacción que creó esa versión.
    \item $\mathrm{MT\text{-}L}(Q_k)$: la mayor marca temporal entre las
    transacciones que leyeron esa versión.
\end{itemize}

Cuando $T_i$ pide $\mathrm{read}(Q)$, el motor elige la versión $Q_k$
con $\mathrm{MT\text{-}E}(Q_k)$ máximo que sea $\le MT(T_i)$: ``la
versión más reciente que ya existía cuando empezó $T_i$''. Para
$\mathrm{write}(Q)$, si $MT(T_i) < \mathrm{MT\text{-}L}(Q_k)$ $T_i$ se
retrocede; si $MT(T_i) = \mathrm{MT\text{-}E}(Q_k)$ se sobrescribe esa
versión; en otro caso se crea una nueva.

\paragraph{Ventaja clave.} Las lecturas \emph{nunca esperan}: hay una
versión apropiada y se devuelve directamente.

\subsection{MVCC en PostgreSQL}

PostgreSQL implementa \textbf{MVCC} (\emph{Multi-Version Concurrency
Control}) y agrega a cada tupla dos columnas de sistema:

\begin{itemize}
    \item \texttt{xmin}: el XID de la transacción que \emph{insertó}
    esa versión.
    \item \texttt{xmax}: el XID de la transacción que la \emph{invalida}
    (por update o delete); $0$ si la versión sigue viva.
\end{itemize}

Cada \texttt{UPDATE} no machaca la tupla: crea una versión nueva con
\texttt{xmin = XID actual} y deja a la versión vieja marcada con
\texttt{xmax = XID actual}. Cuando una transacción quiere leer la fila,
elige la versión que sea visible para su snapshot.

\paragraph{Consecuencia operativa.}
\begin{itemize}
    \item Lecturas no bloquean escrituras, y viceversa.
    \item Dos transacciones concurrentes pueden ver versiones distintas
    de la misma fila simultáneamente (ejercicio 8 del Práctico 6).
    \item Las versiones muertas se eliminan con \texttt{VACUUM} o el
    \texttt{autovacuum}.
\end{itemize}

PostgreSQL combina MVCC con bloqueos explícitos cuando la operación lo
requiere: \texttt{SELECT ... FOR UPDATE}, \texttt{LOCK TABLE}, etc.

\subsection{Modos de bloqueo de tabla en PostgreSQL}

PostgreSQL ofrece varios modos de bloqueo a nivel tabla, en general
adquiridos automáticamente por el comando:

\begin{center}
\begin{tabular}{ll}
\toprule
Modo & Adquirido por \\
\midrule
\texttt{ACCESS SHARE}          & \texttt{SELECT} \\
\texttt{ROW SHARE}             & \texttt{SELECT FOR UPDATE/SHARE} \\
\texttt{ROW EXCLUSIVE}         & \texttt{INSERT}, \texttt{UPDATE}, \texttt{DELETE} \\
\texttt{SHARE UPDATE EXCLUSIVE} & \texttt{VACUUM}, autoanalyze \\
\texttt{SHARE}                 & \texttt{CREATE INDEX} \\
\texttt{SHARE ROW EXCLUSIVE}   & --- (manual) \\
\texttt{EXCLUSIVE}             & --- (manual) \\
\texttt{ACCESS EXCLUSIVE}      & \texttt{ALTER/DROP TABLE}, \texttt{TRUNCATE}, \texttt{REINDEX} \\
\bottomrule
\end{tabular}
\end{center}

Excepto \texttt{ACCESS SHARE}, los modos adquiridos en una transacción se
retienen hasta el final.

% =====================================================
\section{Fenómeno fantasma e índices}
% =====================================================

El control de concurrencia con \emph{granularidad de tupla} (locks por
fila, o marcas por fila) no detecta conflictos \emph{sobre tuplas
todavía no creadas}.

\paragraph{Ejemplo.} $T_1$ ejecuta
\begin{lstlisting}
SELECT min(precio) FROM provee WHERE nart = 2;
\end{lstlisting}
y $T_2$ ejecuta
\begin{lstlisting}
INSERT INTO provee (nart, nprov, precio) VALUES (2, 4, 1.23);
\end{lstlisting}

No hay tupla común; con granularidad de tupla el sistema no ve conflicto.
Sin embargo, $T_1$ podría ver el \texttt{INSERT} de $T_2$ en una segunda
ejecución y ya no sería repetible. Este problema se llama
\textbf{fenómeno fantasma}.

\paragraph{Solución.} Bloquear \emph{a nivel de predicado o de índice}:
un lock sobre \texttt{provee.nart = 2} impide a otras transacciones
insertar filas que cumplan esa condición. Los motores reales lo
implementan vía \emph{predicate locks} o \emph{index range locks}.

% =====================================================
\section{Transacciones en SQL}
% =====================================================

Una transacción SQL queda iniciada implícitamente (al ejecutar la primera
sentencia) o explícitamente (\texttt{BEGIN} / \texttt{START TRANSACTION})
y termina con:

\begin{itemize}
    \item \texttt{COMMIT}: confirma todas las escrituras realizadas.
    \item \texttt{ROLLBACK}: las descarta y restaura el estado anterior.
\end{itemize}

\paragraph{Lo que la norma exige.} El sistema debe garantizar
\emph{secuencialidad} (serializabilidad por conflicto) \emph{y} ausencia
de rollback en cascada. Implementaciones reales suelen relajar la
serializabilidad estricta y dejar al programador elegir el nivel de
aislamiento.

\paragraph{Comandos prácticos.}
\begin{lstlisting}
-- MySQL
SET autocommit = 0;
START TRANSACTION;
... -- operaciones
COMMIT; -- o ROLLBACK

-- PostgreSQL
BEGIN;
... -- operaciones
COMMIT; -- o ROLLBACK
\end{lstlisting}

\paragraph{Constraints diferidas (PostgreSQL).} La cláusula
\texttt{DEFERRABLE INITIALLY DEFERRED} en una FK posterga la verificación
hasta el \texttt{COMMIT}, permitiendo estados intermedios
``inconsistentes'' dentro de la transacción. Es útil cuando dos
\texttt{INSERT} se referencian mutuamente.

% =====================================================
\section{Cómo se conecta con el Práctico 6}
% =====================================================

Los conceptos del Teórico 8 aparecen en el Práctico 6 distribuidos así:

\begin{itemize}
    \item \textbf{Ejercicio 1}: serializabilidad por conflicto + grafo
    de precedencia (sección sobre serializabilidad).
    \item \textbf{Ejercicio 2}: protocolo de bloqueo de dos fases puro
    (sección sobre 2PL).
    \item \textbf{Ejercicio 3}: protocolo de ordenación por marcas
    temporales (sección sobre marcas temporales).
    \item \textbf{Ejercicio 4}: atomicidad y \texttt{ROLLBACK} en
    MySQL (sección sobre transacciones en SQL).
    \item \textbf{Ejercicio 5}: nivel \texttt{SERIALIZABLE} y bloqueos
    en MySQL (sección sobre niveles de aislamiento y bloqueos).
    \item \textbf{Ejercicio 6}: lectura no repetible y niveles
    \texttt{READ COMMITTED} / \texttt{REPEATABLE READ} (sección sobre
    niveles de aislamiento).
    \item \textbf{Ejercicio 7}: constraints deferibles en PostgreSQL
    (sección sobre transacciones en SQL).
    \item \textbf{Ejercicio 8}: MVCC, \texttt{xmin}, \texttt{xmax}
    (sección sobre esquemas multiversión).
\end{itemize}

% =====================================================
\section{Conclusión}
% =====================================================

Las ideas centrales a retener:

\begin{itemize}
    \item Una transacción es la unidad lógica del trabajo concurrente; el
    motor garantiza para ella las propiedades ACID.
    \item El control de concurrencia tiene que asegurar
    \emph{aislamiento}: el resultado de la ejecución intercalada debe
    ser equivalente a alguna ejecución serial.
    \item El test del grafo de precedencia es la herramienta formal para
    decidir si un schedule es serializable por conflicto.
    \item Hay dos grandes familias de protocolos: basados en bloqueos
    (2PL y variantes) y basados en marcas temporales. Los motores
    modernos suelen combinar ambos y agregar multiversión (MVCC).
    \item Los niveles de aislamiento (\texttt{READ UNCOMMITTED},
    \texttt{READ COMMITTED}, \texttt{REPEATABLE READ},
    \texttt{SERIALIZABLE}) ofrecen un trade-off explícito entre
    consistencia y concurrencia, regulando qué fenómenos pueden aparecer.
    \item Recuperabilidad y ausencia de rollback en cascada son
    propiedades adicionales que la planificación debe cumplir para que
    el sistema pueda revertir transacciones de manera limpia.
\end{itemize}

\end{document}
